% Part: sets-functions-relations
% Chapter: size-of-sets-complete

\documentclass[../../../include/open-logic-chapter]{subfiles}

\begin{document}

\olchapter{sfr}{siz}{De grootte van verzamelingen}

\begin{editorial}
Dit hoofdstuk behandelt opsommingen, aftelbaarheid en
overaftelbaarheid. Verscheidene paragrafen bestaan in twee versies: een
meer elementaire versie die opsommingen opvat als lijsten of als
surjecties vanuit $\PosInt$, en een abstractere versie die opsommingen
definieert als bijecties met $\Nat$.
\end{editorial}

\olimport{introduction}

\olimport{enumerability}

\olimport{zig-zag}

\olimport{pairing}

\olimport{pairing-alt}

\olimport{non-enumerability}

\olimport{reduction}

\olimport{equinumerous-sets}

\olimport{comparing-size}

\olimport{schroder-bernstein}

\begin{editorial}
De volgende paragrafen, \olref[sfr][siz][enm-alt]{sec},
\olref[sfr][siz][nen-alt]{sec} en \olref[sfr][siz][red-alt]{sec}, zijn
alternatieve versies van \olref[sfr][siz][enm]{sec},
\olref[sfr][siz][nen]{sec} en \olref[sfr][siz][red]{sec}. Tim Button
schreef ze voor zijn tekst Open Set Theory. Ze zijn iets gevorderder
en gebruiken een andere definitie van opsombaarheid die beter past in
de verzamelingenleer: een bijectie met $\Nat$ of met een beginsegment,
in plaats van de eis dat iets in een lijst kan worden gezet of het
beeld is van een surjectieve functie vanuit $\PosInt$.
\end{editorial}

\olimport{enumerability-alt}
\olimport{non-enumerability-alt}
\olimport{reduction-alt}

\OLEndChapterHook
\end{document}
